% !TEX root = ../main.tex
\section{Logically coherent \summs\ as probabilities}
\label{sec:coh_prob}

This section shows that
a \summ \ is logically coherent if and only if
it is a coherent representation of uncertainty \citep{Finetti1931}, that is,
it admits a probability representation.
While \cref{sec:coherent_standard} 
discusses the special case of 
\stdsumms, \cref{sec:coherent_agnostic} discusses
thernary \summs.

\subsection{\Stdsumms}
\label{sec:coherent_standard}

From a practical perspective,
a \stdsumm \ divides hypotheses
into those that 
are summarized as true $(\phi(H|x) = 1)$
and as false $(\phi(H|x) = 0)$.
One might ask whether
the latter can be interpreted as
having probability $0$ and
the former as having probability $1$.
\Cref{cor:std_prob} shows that
this is the case if and only if
$\phi$ is at least 
a $\aleph_0$-coherent \summ.

\begin{theorem}
 \label{cor:std_prob}
 A \stdsumm, $\phi$, is $\aleph_0$-coherent
 if and only if $\phi$ is
 a $\{0,1\}$-probability.
\end{theorem}

However, what is the difference between
$\aleph_0$-coherence and
other types of $\kappa$-coherence? 
The following result shows that,
under weak conditions,
$\kappa$-coherence can 
divided into only two types:

\begin{theorem}
 \label{thm:coherent_std_1}
 If $|\Theta|$ is not a measurable cardinal
 \citep{Ulam1930} and $\phi$ is $\aleph_1$-coherent, then $\phi$ is completely-coherent, that is, $\phi$ is a completely-additive $\{0,1\}$-probability.
\end{theorem}

In other words,
if $|\Theta|$ is not a measurable cardinal, then
$\aleph_0$-coherent \stdsumms\ are
divided into two types.
Either $\phi$ is not $\aleph_1$-coherent and
is an extreme-valued 
strictly finitely-additive probability, or
$\phi$ is $\aleph_1$-coherent and
an extreme-valued
completely-additive probability.

The next subsection shows that
a similar conclusion exists
for logically-coherent 
thernary \summs.

\subsection{General \summs}
\label{sec:coherent_agnostic}

A three-valued \summ\ divides 
hypotheses among
the ones that are believed to be true
$(\phi(H|x) = 1)$, 
to be false
$(\phi(H|x) = 0)$,
and the ones that remain unknown
$(\phi(H|x) = \half)$.
From a practical perspective, 
one might interpret 
true hypotheses as ``highly likely'',
false hypotheses as ``highly unlikely'', and
undecided hypotheses as somewhere in-between.
\Cref{def:rep} is a necessary condition for
such an interpretation:

\begin{definition}
 \label{def:rep}
 A nonstandard probability, $\P_x$, 
 represents a \summ, $\phi$, if
 for every $x \in \sX$ and $\alpha \in (0,1)$:
 \begin{itemize}
  \item If $\phi(H_A^1|x) = \phi(H_A^2|x) = 1$,
  then $\alpha\P_x(H_A^1) < \P_x(H_A^2)$,
  \item If $\phi(H_A|x) = 1$ and 
  $\phi(H_U|x) = \half$, then
  $\P_x(H_A^c) < 
  \alpha \min(\P_x(H_U^c), \P_x(H_U))$,
  \item If $\phi(H_A|x) = 1$ and 
  $\phi(H_U|x) = 0$, then
  $\P_x(H_A^c) < 
  \alpha \P_x(H_U^c)$ and
  $\P_x(H_R) < \alpha \P_x(H_U)$,
  \item  If $\phi(H_R|x) = 0$ and
  $\phi(H_U|x) = \half$, then
  $\P_x(H_R) 
  < \alpha \min(\P_x(H_U),\P_x(H_U^c))$.
 \end{itemize}
\end{definition}

In other words, if $\P_x$ represents $\phi$,
then the probability of 
every undecided hypothesis and its complement is
in-between that of every hypothesis
believed to be true false and 
believed to be true.
Furthermore, the probability of every
hypothesis belived to be false is 
arbitrarily closer to $0$
then every other hypothesis 
not believed to be false.
Similarly,
the probability of every hypothesis
belived to be true
is arbitrarily closer to $1$ then 
every other not believed to be true.

The existence of such a probability representation
is equivalent to $\aleph_0$-coherence:

\begin{theorem}
 \label{thm:coherence-nonstd}
 A proper \summ, $\phi$, is 
 $\aleph_0$-coherent
 if and only if
 there exists a non-standard probability
 that represents $\phi$. 
\end{theorem}

\cref{thm:coherence-nonstd} shows that,
$\phi$ is $\aleph_0$-coherent
if and only if
it agrees qualitatively with
a nonstandard probability, $\P_x$.
That is, the standard part of $\P_x$ is $0$
for every false  hypothesis and is $1$
for every true hypothesis.
Furthermore, the probability of 
every undecided hypothesis is
higher than that of every false hypothesis and
lower than that of every true hypothesis.

Particular types of $\aleph_0$-coherent \summs \ are region-based \summs, as described below:

\begin{definition}[Region-based \summ]
 \label{def:region-based}
 $\phi$ is a region-based \summ \ if
 there exists $R_x \subset \Theta$ such that
 \begin{align*}
  \phi(H|x) &=
  \begin{cases}
   1 & \text{, if } R_x \subseteq H, \\
   0 & \text{, if } R_x \cap H = \emptyset, \\
   \half & \text{, otherwise.}
  \end{cases}
 \end{align*}
\end{definition}

Region-based \summs \ are illustrated
in \cref{fig::region} and
admit an intuitive interpretation. 
Assume that $R_x$ are
plausible points for $\theta$ 
after $x$ is observed.
If all plausible points fall in $H$,
then one believes that $H$ is true.
Similarly, if all plausible points
fall outside of $H$, then
one believes $H$ to be false.
Otherwise, one remains undecided.
Furthermore, a particular representation
of a region-based \summ \ assigns
probability $1$ to $R_x$ and
an infinitesimal value to each element of $R_x$.

\begin{lemma}
 \label{thm:region-is-coherent}
 If $\phi$ is a region-based test,
 then $\phi$ is completely-coherent.
\end{lemma}

\begin{figure}
  \center
  \gfbstfig
  \mbox{} \vspace{-2mm} \mbox{} \\ 
  \caption{$\phi$ is
  a region-based \summ \ for $H$.}
  \label{fig::region}
\end{figure}

\begin{example}
\label{sec:scheffe}
Both Scheff\`{e}'s test for contrasts \citep{Scheffe1953} and the standard ANOVA $F$-test can be combined into a single region-based \summ. Consider the one-way analysis of variance layout in which, for groups $i \in \{1,\dots,k\}$,
\begin{align*}
 Y_{ij} &= \mu_i + \epsilon_{ij},
 \quad j \in \{1, \dots, n_i\},
 \quad \epsilon_{ij}
 \overset{iid}{\sim} N(0,\sigma^2).
\end{align*}
The sample size is $N = \sum_{i} n_i$ and the unknown quantity is $\theta = (\mu,\sigma^2)$, with $\mu = (\mu_1,\dots,\mu_k) \in \Re^k$.
Common frequentist estimators for these parameters are the group averages, $\hat{\mu}_i := \bar{Y}_i = n_i^{-1}\sum_{j} Y_{ij}$, and the pooled variance estimate, $\hat{\sigma}^2 = (N-k)^{-1} \sum_{i,j}(Y_{ij} - \bar{Y}_i)^2$.

Let $\mathcal{C} = \{c \in \Re^k: \sum_{i} c_i = 0\}$ be the set of contrasts. A contrast $c \in \mathcal{C}$ specifies the hypothesis, $H_c = \{\mu \in \Re^k: c \cdot \mu = 0\}$, that the corresponding linear combination of the group means vanishes. Scheff\`{e}'s method provides a statistical hypothesis test for each $H_c$. Specifically, $H_c$ is rejected if
\begin{align*}
 \frac{\left|\sum_{i} c_i \bar{Y}_i\right|^2}{\hat{\sigma}^2 \sum_{i} c_i^2/n_i}
 &> (k-1) F_{\alpha; \, k-1, \, N-k},
\end{align*} 
where $F_{\alpha;\,k-1,\,N-k}$ is the $1-\alpha$ quantile of the $F$ distribution with $k-1$ and $N-k$ degrees of freedom. 

Also, the omnibus null hypothesis is that all group means coincide, that is,
\begin{align*}
 H_0 &= \{\mu \in \Re^k: \mu_1 = \dots = \mu_k\}
     = \cap_{c \in \mathcal{C}} H_c.
\end{align*}
The standard ANOVA $F$-test rejects $H_0$ if
\begin{align*}
 \frac{\sum_i n_i
(\bar{Y}_i - \bar{Y})^2}{\hat{\sigma}^2} &> (k-1) F_{\alpha; \, k-1, \, N-k}.
\end{align*}

It is possible to describe both Scheff\`{e}'s method and the $F$-test as a single \summ, $\phi$. For this purpose, let $R_x$ be the standard confidence region for $\mu$:
\begin{align*}
 R_x &= \left\{\mu \in \Re^k: \frac{\sum_{i} n_i \left[(\bar{Y}_i - \mu_i) - (\bar{Y} - \bar{\mu})\right]^2}{\hat{\sigma}^2}
 \leq (k-1) F_{\alpha; \, k-1, \, N-k}\right\}.
\end{align*}

Consider the \summ, $\phi$, defined as follows:
\begin{align*}
 \phi(H|x) &=
 \begin{cases}
  1 & \text{, if } R_x \subseteq H, \\
  0 & \text{, if } R_x \cap H = \emptyset, \\
  \half & \text{, otherwise.}
 \end{cases}
\end{align*}
$\phi$ describes both Scheff\`{e}'s method and the $F$-test. For every $c \in \mathcal{C}$, $\phi(H_c|x) = 0$ if and only if Scheff\`{e}'s method rejects $H_c$. Similarly, $\phi(H_0|x) = 0$ if and only if $H_0$ is rejected by the $F$-test. Furthermore, it follows from \cref{thm:region-is-coherent} that $\phi$ is completely-coherent. In particular, if there exists $H_c$ such that $\phi(H_c|x) = 0$, then $\phi(H_0|x) = 0$ \citep{Gabriel1969}.
\end{example}

Under weak conditions, a \summ \ is
$\aleph_1$-coherent if and only if
it is a region-based \summ.

\begin{definition}
 For a topology over $\Theta$, $\tau$, 
 and $H \subset \Theta$, let
 $\text{Cl}_\tau(H)$ denote the
 closure of $H$ with respect to $\tau$. 
 A \summ, $\phi$, is smooth with respect to
 a topology, $\tau$, if
 $\phi(H|x) = 0$ implies that
 $\phi(\text{Cl}_\tau(H)|x) = 0$.
\end{definition}

\begin{theorem}
 \label{thm:second-countable}
 Assume that there exists 
 a second-countable topology, $\tau$, 
 such that $\phi$ is smooth
 with respect to $\tau$ and
 $\tau \subset \sigma(\Theta)$.
 The \summ, $\phi$, is $\aleph_1$-coherent
 if and only if
 it is a region-based \summ.
\end{theorem}

\Cref{thm:second-countable} shows that,
similarly to \stdsumms,
$\aleph_0$-coherent 
\summs \ are also 
divided into two types.
If a \summ \ is not $\aleph_1$-coherent, then
it admits a probability representation and
is not a region-based \summ.
If the \summ \ is $\aleph_1$-coherent, then
it is a region-based \summ.

\red{Not sure how to connect the following examples to the text.}

\begin{example}
 Let $\Theta \subset \Re^d$, and $\phix(H) = \I(\P(\theta \in H > k|x))$, for some $k \in (0,1)$. If for every $x$, $\P(\theta \leq t|X=x)$ is a continuous function of $t$, then $\phi$ is not smooth with respect to the usual topology over $\Re^d$. For instance, for every $x$, $\phi(\mathbb{Q}^d|x) = 0$ and $\phi(\Re^d|x) = 1$.
\end{example}

\begin{example}
 Let $\Theta \subset \Re^d$ and
 $\mu$ be the Lebesgue measure over $\Theta$.
 The log-likelihood of $\theta \in \Theta$
 based on $x \in \sX$ is defined as
  $l_{x}(\theta)
 := \log\left(\frac{d\P_{\theta}(x)}{d\mu}\right)$.
 Similarly, the log-likelihood statistic for
 an hypothesis, $H$, is
 $\Lambda_{H}(x) :=
 -2\left(\sup_{\theta \in H}l_x(\theta) 
 - \sup_{\theta \in \Theta}l_x(\theta)\right)$.
 The standard log-likehood \summ \ is defined as
 $\phi(H|x) = \I(\Lambda_{H}(x) > k_{H})$,
 where $k_H$ is chosen so that
 $\sup_{\theta \in H} \E_{\theta \in H}[\phi(H)] = \alpha$.
 If the log-likelihood is continuous with respect to $\tau$,
 then $\phi$ is smooth 
 with respect to $\tau$.
\end{example}


